\documentclass[crop=false]{standalone}
% --- ПРЕАМБУЛА ДЛЯ ПОДДЕРЖКИ РУССКОГО ЯЗЫКА И МАТЕМАТИКИ ---
\usepackage[T2A]{fontenc}
\usepackage[utf8]{inputenc}
\usepackage[russian]{babel}
\usepackage{amsmath}
\usepackage{geometry}
\usepackage{amssymb}
\usepackage{graphicx}
\usepackage{float}
\usepackage{import}
\geometry{a4paper, top=2cm, bottom=2cm, left=2cm, right=2cm}
\newcommand{\sidecomment}[1]{\marginpar{\raggedright\scriptsize\itshape #1}}
\newcommand{\iu}{\mathrm{i}}

\newcommand{\Def}{%
    \text{Def}%
    \hspace{-0.85em}% Сдвиг назад ровно на середину (подберите под шрифт)
    \raisebox{0.1ex}{/}% Черта
    \hspace{1em}% Возврат к концу слова + небольшой отступ
}

\renewcommand{\Re}{\operatorname{Re}}
\renewcommand{\Im}{\operatorname{Im}}

\ifstandalone
    \newcommand{\lecturebreak}{} % В отдельном файле лекции разрыв не нужен
\else
    \newcommand{\lecturebreak}{\clearpage} % В полной книге - начинаем с новой страницы
\fi

\newcounter{lecturenumber} % Создаем счетчик для нумерации лекций
\newcommand{\lecture}[2]{
    \lecturebreak
    \refstepcounter{lecturenumber} % Увеличиваем счетчик на 1
    \par\vspace{2em}\noindent % Отступ сверху и убираем абзацный отступ
    {\Large\bfseries % Делаем заголовок крупным и жирным
    Лекция \thelecturenumber: #1 % Выводим "Лекция N: Тема"
    \hfill % Растягиваем пробел, чтобы дата была справа
    \large\textit{#2}} % Выводим дату курсивом
    \par\vspace{1em} % Небольшой отступ снизу
    \addcontentsline{toc}{section}{Лекция \thelecturenumber: #1} % Добавляем в оглавление
    \par
}

\newcommand{\TwoCols}[2]{%
  \par\noindent % Начать с новой строки без отступа
  \begin{minipage}[t]{0.48\textwidth}
    #1
  \end{minipage}% <-- Важный '%' для удаления пробела
  \hfill % Растягивающийся пробел между колонками
  \begin{minipage}[t]{0.48\textwidth}
    #2
  \end{minipage}% <-- Важный '%' для удаления пробела
  \par%
}

\pagestyle{plain} % Включаем нумерацию страниц\
\begin{document}
\lecture{Кинематические уравнения для Кватернионов}{24.02}
\subsection{Замечания к предыдущей лекции}
\keypoint вместо $\Lambda$ можно брать не единичные кватернионы

\keypoint \[ 
\begin{aligned}
\| q_1 - q_2 \| \stackrel{?}{=} \| \psi_\Lambda(q_1) - \psi_\Lambda(q_2) &= \| \Lambda \circ q_1 \circ  \Lambda^{-1} - \Lambda \circ q_2 \circ \Lambda^{-1} \| \\ 
&= \| \Lambda \circ (q_1 - q_2) \circ \Lambda^{-1} \| = \Lambda \circ (q_1 - q_2) \circ \Lambda^{-1} \circ \overline{\Lambda \circ (q_1 - q_2) \circ \Lambda^{-1}} \\ 
&= \Lambda \circ (q_1 - q_2) \circ \Lambda^{-1} \circ \bar \Lambda^{-1} \circ \overline {(q_1 - q_2)} \circ \bar \Lambda = \frac{1}{\| \Lambda \|} \Lambda \circ (q_1 - q_2) \circ \overline{(q_1 - q_2)} \circ \bar \Lambda \\ 
&= \frac{1}{\| \Lambda \|} \| q_1 - q_2 \| \Lambda \circ \bar \Lambda = \| q_1 -  q_2 \|
\end{aligned}
\]
Доказали, что $\psi$ сохраняет норму вектора

\subsection{Кинематические уравнения для кватернионов}
\keypoint Кватернион поворота на малый угол $\Delta \varphi$
\[\Lambda = \cos \frac{\Delta \varphi}{2} + \sin \frac{\Delta \varphi}{2} \bar u = 1 + \frac{\Delta \varphi}{2} \bar u + O(\Delta \varphi^2)\]

\[\boxed{ \vec \omega = \lim_{\Delta t \to 0} \frac{\Delta \varphi}{\Delta t} \vec u(t , \Delta t)} \]
%fixme
\[
\begin{aligned}
\dot \Lambda = \lim_{\Delta t \to 0} \frac{\Lambda(t + \Delta t) - \Lambda(t)}{\Delta t} &\stackrel{\text{АТЗ}}{=} \lim_{\Delta t \to 0} \frac{\Lambda(\Delta t) \circ \Lambda(t) - \Lambda(t)}{\Delta t} = \left(\lim_{\Delta t \to 0} \frac{\Lambda(\Delta t) -1 }{\Delta t}\right)\circ \Lambda(t) \\ 
&= \left( \lim_{\Delta t \to 0} \frac{\Delta \varphi \bar u}{2 \Delta t}\right) \circ \Lambda(t) = \frac{1}{2} \bar \omega \circ \Lambda(t)
\end{aligned}
\]
Активная точка зрения:
\[
\boxed{
\dot \Lambda = \frac{1}{2} \vec \omega  \circ \Lambda
}
\]
Упраженние: записать в пассивной

\section{Кинематика систем со связями}
$N$ Точек тела

$\bar r_1, \dots , \bar r_N$ — законы движения, обозначается $\bar r_k(t)$

$\bar v_1, \dots , \bar v_N$ — скорости
\Def Механическая связь — соотношение
\[f (\bar r_1 , \dots, \bar r_n , \bar v_1, \dots , \bar v_N, t) = 0\]
Удерживающая связь — равно, неудерживающая — меньше равно

\begin{table}[H]
    \centering
    \begin{tabular}{c|c|c}
        & Геом & Дифф \\
        Стац & $f(\bar r_1, \dots , \bar r_N) = 0$ & $f(\bar r_1, \dots , \bar r_N, \bar v_1, \dots , \bar v_N) = 0 \qquad \sum^N_{k=1} \bar a_k (\bar r_1, \dots, \bar r_N) \bar v_k = 0$ \\ % примеры диф. сввяей, которые встречаются(линейные по скоростям)
        Нест & $f(\bar r_1, \dots, \bar r_N, t) = 0$ & $f(\bar r_1, \dots ,\bar r_N , \bar v_1, \dots , \bar v_N, t) = 0 \qquad \sum^N_{k=1} \bar a_k (\bar r_1, \dots, \bar r_N, t) \bar v_k + a_0 (\bar r_1, \dots ,\bar r_N, t) = 0$
    \end{tabular}
    \caption{}
\end{table}

\[f(\bar r_k , t) = 0 \qquad \dot f(\bar r_k , t) = \sum^N_{j=1} \frac{\partial f(\bar r_k, t)}{\partial \bar r_j} \frac{d \bar r_j}{dt} + \frac{\partial f(\bar r_k, t)}{\partial t} = 0\]
\[\sum^N_{j=1} \frac{\partial f}{\partial \bar r_J} \bar v_j + \frac{\partial f}{\partial t} = 0 \]
\Def Дифференцируемая связь, которая эквивалентна некоторой геометрической называется интегрируемой, иначе неинтегрируемой

\keypoint Движение по окружности, Связи
\[ x^2 + y^2 - R^2 = 0 \qquad (\bar r , \bar v) = 0 \qquad xv_x + y v_y = 0\]
эквивалентны
\begin{proof}
    \[x^2 + y^2 - R^2 = 0 \qquad 2 x \dot x + 2 y \dot y = 0\]
\end{proof}
\noindent \keypoint Пример неинтегрируемой связи
% fixme картинка в телефоне
Вторая точка двигается по направлению к первой
\[\bar v_1 = \alpha (\bar r_0 - \bar r_1) \qquad v_{1x} = \alpha (x_0 - x_1) \qquad v_{1y} = \alpha(y_0 - y_1) \qquad \frac{v_{1x}}{v_{1y}} = \frac{x - x_0}{y_0 - y_1}\]
\begin{proof}
    \[v_{1x} (y_0 - y_1) - v_{1y}(x_0 - x_1) = 0\]
    От противного:
    существует $f$:
    \[\frac{\partial f}{\partial \bar r_0} \bar v_0 + \frac{\partial f}{\partial \bar r_1} \bar v_1 = v_{1x} (y_0 -y_1) - v_{1y}(x_0 -x_1) = 0\]
    \[\frac{\partial f}{\partial \bar r_0} = 0 \qquad \frac{\partial f}{\partial t} = 0\]
    \[\frac{\partial f}{\partial \bar r_1} = \begin{pmatrix}
        y_0 - y_1 \\ -(x_0 - x_1)
    \end{pmatrix} \iff \begin{matrix}
        \frac{\partial f}{\partial x_1} = y_0 - y_1 \\ \\ \frac{\partial f}{\partial y_1} = x_1 - x_0
    \end{matrix} \qquad \begin{matrix} \frac{\partial^2 f}{\partial x_1 \partial y_1} = -1 \\ \\  \frac{\partial^2 f}{\partial y_1 \partial x_1} = 1 \end{matrix}
    \]
    Противоречие с тем, что у непрерывной функции перекрёстные производные разные
\end{proof}
% fixme есть в телефоне записи дениса, разобраться

$N: \bar r_1, \dots , \bar r_N$
\[f_P(\bar r_1, \dots, \bar r_N , t) = 0 \qquad p = 1, \dots , P\]
\[\sum_{k=1}^N \bar a_{ks}(\bar r_1, \dots, \bar r_N, t) \bar v_K + a_s (\bar r_1, \dots , \bar r_N, t) \qquad s = 1, \dots , S\]
Далее будем считать, что если связь дифференцируемая, то она неинтегрируемая, иначе сопоставляем ей геометрическую
\Def Голономная система — система, в которой вся связи только геометрические, иначе Неголономная
\Def Склерономная система — система, в которой все связи стационарны, иначе реаномная
\begin{table}
    \begin{tabular}{c|c|c}
      & Голономная & Неголономная \\
      Склерономная & $\checkmark$  & \\
      Реаномная & $\checkmark$  &   
    \end{tabular}
\end{table}
Упражение: согнуть ложку, чтобы получился кельтский камень

По часовой будет крутиться, против часовой крутиться не будет

Такая система неголономна
%fixme есть картинка в телефоне
\end{document}
